
\documentclass{article}
\usepackage{colortbl}
\usepackage{makecell}
\usepackage{multirow}
\usepackage{supertabular}

\begin{document}

\newcounter{utterance}

\centering \large Interaction Transcript for game `cladder', experiment `full\_v1.5\_default', episode 4040 with reflexion{-}Qwen3.5{-}4B.
\vspace{24pt}

{ \footnotesize  \setcounter{utterance}{1}
\setlength{\tabcolsep}{0pt}
\begin{supertabular}{c@{$\;$}|p{.15\linewidth}@{}p{.15\linewidth}p{.15\linewidth}p{.15\linewidth}p{.15\linewidth}p{.15\linewidth}}
    \# & \multicolumn{2}{c}{Player} && \multicolumn{2}{c}{Game Master} \\
    \hline

    \theutterance \stepcounter{utterance}  
    & & & \multicolumn{4}{p{0.6\linewidth}}{
        \cellcolor[rgb]{0.9,0.9,0.9}{
            \makecell[{{p{\linewidth}}}]{
                \texttt{\tiny{[P1$\langle$GM]}}
                \texttt{You have attempted this same task before and did not succeed. These are the plans you wrote after those attempts.} \\
\texttt{Trial \#1: New plan: I failed because I misinterpreted the specific question asked. The prompt asked "Will smoking decrease the chance...", and I initially answered "No" (likely confusing the question with a negative phrasing or getting lost in the causal structure details before checking the final conclusion). Although I later corrected myself in the text, the final output was incomplete and did not provide a definitive "Yes" or "No" as the very first word, nor did it fully conclude the reasoning. The strategy must be to immediately answer "Yes" or "No" based strictly on the comparison of the given probabilities (\$37\textbackslash{}\% < 77\textbackslash{}\%\$), then provide the causal reasoning that confirms this numerical difference represents a decrease. I will ensure the final answer is a complete, self{-}contained statement starting with the correct boolean response.} \\
\\ 
\\ 
\texttt{You are an expert in causal inference. The following question is not a typical commonsense query, but rather a meticulously designed question created by a professor specializing in causal inference, intended to assess the students' mastery of the course content.} \\
\texttt{Imagine a self{-}contained, hypothetical world with only the following conditions, and without any unmentioned factors or causal relationships: Smoking has a direct effect on effort and lactose intolerance. Effort has a direct effect on lactose intolerance. For nonsmokers, the probability of being lactose intolerant is 77\%. For smokers, the probability of being lactose intolerant is 37\%. Will smoking decrease the chance of being lactose intolerant?} \\
\texttt{Start your answer with 'Yes' or 'No', followed by additional reasoning or evidence to support your explanation.} \\
            }
        }
    }
     \\ \\

    \theutterance \stepcounter{utterance}  
    & \multicolumn{4}{p{0.6\linewidth}}{
        \cellcolor[rgb]{0.561,0.737,0.561}{
            \makecell[{{p{\linewidth}}}]{
                \texttt{\tiny{[P1$\rangle$GM]}}
                \texttt{Yes. In this hypothetical world, smoking directly affects effort, which in turn directly affects lactose intolerance. We are given the marginal probabilities: the probability of being lactose intolerant for nonsmokers is 77\%, and for smokers, it is 37\%. Since 37\% is strictly less than 77\%, the presence of smoking is associated with a lower probability of the outcome. Because the prompt explicitly states that there are no unmentioned factors or causal relationships, this observed difference in probabilities must be attributed to the causal pathway where smoking influences effort, which subsequently reduces the likelihood of lactose intolerance. Therefore, smoking decreases the chance of being lactose intolerant.} \\
            }
        }
    }
    & & \\ \\

    \theutterance \stepcounter{utterance}  
    & & & \multicolumn{2}{p{0.3\linewidth}}{
        \cellcolor[rgb]{0.9,0.9,0.9}{
            \makecell[{{p{\linewidth}}}]{
                \texttt{\tiny{[GM$|$GM]}}
                \texttt{Yes} \\
            }
        }
    }
    & & \\ \\

    \theutterance \stepcounter{utterance}  
    & & & \multicolumn{2}{p{0.3\linewidth}}{
        \cellcolor[rgb]{0.9,0.9,0.9}{
            \makecell[{{p{\linewidth}}}]{
                \texttt{\tiny{[GM$|$GM]}}
                \texttt{yes} \\
            }
        }
    }
    & & \\ \\

    \theutterance \stepcounter{utterance}  
    & & & \multicolumn{2}{p{0.3\linewidth}}{
        \cellcolor[rgb]{0.9,0.9,0.9}{
            \makecell[{{p{\linewidth}}}]{
                \texttt{\tiny{[GM$|$GM]}}
                \texttt{game\_result = WIN} \\
            }
        }
    }
    & & \\ \\

\end{supertabular}
}

\end{document}
